\documentclass[11pt]{article}
\usepackage{amsmath,amssymb,amsthm}
\usepackage{geometry}
\usepackage{enumitem}
\geometry{margin=1in}

\title{\textbf{CSE400 -- Probability}\\
Lecture 6 (L6)}
\author{Instructor: Dhaval Patel, PhD}
\date{January 22, 2026}

\begin{document}
\maketitle

\section{Introduction and Scope of the Scribe}

These notes faithfully reconstruct \textbf{Lecture 6} as delivered in class.
The goal is to serve as a \textbf{standalone exam reference}, reproducing
definitions, notation, derivations, proofs, and worked examples exactly as
presented in the lecture slides and board work.
Only material explicitly covered in the lecture is included.

\section{Random Variables}

\subsection{Definition}

A \textbf{random variable (RV)} $(X)$ on a sample space $(\Omega)$ is a function
\[
X : \Omega \to \mathbb{R}
\]
that assigns to each sample point $(\omega \in \Omega)$ a real number
$(X(\omega))$.

Until further notice, attention is restricted to \textbf{discrete random
variables}, i.e., random variables whose range is finite or countably infinite.

Although $(X)$ maps $(\Omega)$ to $(\mathbb{R})$, the actual set of values
\[
\{X(\omega) : \omega \in \Omega\}
\]
is a discrete subset of $(\mathbb{R})$.

\subsection{Visualization of a Distribution}

The distribution of a discrete random variable can be visualized using a
\textbf{bar diagram}:
\begin{itemize}
  \item The x-axis represents the values the random variable can take.
  \item The height of the bar at value $(a)$ is $\Pr[X = a]$.
\end{itemize}

Each probability is computed by identifying the corresponding event in the
sample space.

\section{Examples of Random Variables}

\subsection*{Example 3.1: Tossing Three Fair Coins}

\textbf{Experiment:} Toss 3 fair coins.

\textbf{Random Variable:} Let $(Y)$ denote the number of heads observed.

\textbf{Possible values:}
\[
Y \in \{0,1,2,3\}.
\]

\textbf{Probabilities:}
\[
\begin{aligned}
\Pr(Y=0) &= \Pr((T,T,T)) = \tfrac{1}{8},\\
\Pr(Y=1) &= \Pr((H,T,T),(T,H,T),(T,T,H)) = \tfrac{3}{8},\\
\Pr(Y=2) &= \Pr((H,H,T),(H,T,H),(T,H,H)) = \tfrac{3}{8},\\
\Pr(Y=3) &= \Pr((H,H,H)) = \tfrac{1}{8}.
\end{aligned}
\]

\section{Probability Mass Function (PMF)}

\subsection{Definition}

A random variable that can take on at most a countable number of possible values
is said to be \textbf{discrete}.

Let $(X)$ be a discrete random variable with range
\[
R_X = \{x_1, x_2, x_3, \ldots\}.
\]

The function
\[
p_X(x_k) = \Pr(X = x_k), \quad k = 1,2,3,\ldots
\]
is called the \textbf{Probability Mass Function (PMF)} of $(X)$.

Since $(X)$ must take exactly one of these values,
\[
\sum_k p_X(x_k) = 1.
\]

\subsection{PMF Example}

\textbf{Example 4.1:} Two independent tosses of a fair coin.

The PMF was constructed by enumerating outcomes and assigning probabilities
accordingly.

\subsection{PMF with Unknown Constant}

\textbf{Example 4.2:}

The PMF of a random variable $(X)$ is given by
\[
p(i) = C e^{-\lambda i}, \quad i = 0,1,2,\ldots
\]
where $\lambda > 0$.

\textbf{Step 1: Find $(C)$}
\[
\sum_{i=0}^{\infty} C e^{-\lambda i} = 1.
\]
Using the geometric series,
\[
\sum_{i=0}^{\infty} e^{-\lambda i} = \frac{1}{1 - e^{-\lambda}},
\]
we get
\[
C = 1 - e^{-\lambda}.
\]

\textbf{Step 2: Compute probabilities}
\[
\Pr(X>2) = 1 - \Pr(X=0) - \Pr(X=1) - \Pr(X=2).
\]

\section{Bayes' Theorem}

\subsection{Statement (Recap)}

Using
\[
\Pr(A \cap B_i) = \Pr(B_i | A)\Pr(A),
\]
we obtain the \textbf{Bayes Formula}:
\[
\Pr(B_i | A) =
\frac{\Pr(A | B_i)\Pr(B_i)}
{\sum_j \Pr(A | B_j)\Pr(B_j)}.
\]

\begin{itemize}
  \item $\Pr(B_i)$: \emph{a priori probability}
  \item $\Pr(B_i|A)$: \emph{posteriori probability}
\end{itemize}

\subsection{Example 5.1: Auditorium with 30 Rows}

\begin{itemize}
  \item Row 1 has 11 seats, Row 2 has 12 seats, $\ldots$, Row 30 has 40 seats.
  \item A row is chosen uniformly at random.
  \item A seat is chosen uniformly at random within the selected row.
\end{itemize}

\textbf{Tasks:}
\begin{enumerate}
  \item Compute $\Pr(\text{Seat 15} \mid \text{Row 20})$.
  \item Compute $\Pr(\text{Row 20} \mid \text{Seat 15})$.
\end{enumerate}

These were computed using Bayes' theorem exactly as shown in lecture.

\subsection{Example 5.2: Communication System}

A binary communication system transmits bits $0$ or $1$.

\textbf{Conditional probabilities:}
\[
\begin{aligned}
\Pr(0 \text{ received }|0 \text{ transmitted}) &= 0.95,\\
\Pr(1 \text{ received }|0 \text{ transmitted}) &= 0.05,\\
\Pr(0 \text{ received }|1 \text{ transmitted}) &= 0.10,\\
\Pr(1 \text{ received }|1 \text{ transmitted}) &= 0.90.
\end{aligned}
\]

\textbf{Questions addressed in lecture:}
\begin{enumerate}
  \item Find $\Pr(0 \text{ received})$ and $\Pr(1 \text{ received})$.
  \item Given a received bit, find the probability of the transmitted bit.
  \item Find the overall probability of error.
\end{enumerate}

All results were derived step by step using Bayes' theorem.

\section{Independent Events}

\subsection{Definition}

Two events $(A)$ and $(B)$ are \textbf{independent} if
\[
\Pr(A|B) = \Pr(A)
\quad \text{and} \quad
\Pr(B|A) = \Pr(B).
\]

Equivalently,
\[
\Pr(A \cap B) = \Pr(A)\Pr(B).
\]

\subsection{Example: Modified Auditorium}

If each row has the same number of seats:
\begin{itemize}
  \item Let $(A)$ = Row 20 selected
  \item Let $(B)$ = Seat 15 selected
\end{itemize}

\textbf{Conclusion:} Event $(B)$ provides no additional information about $(A)$.
Hence, $(A)$ and $(B)$ are independent.

\subsection{Mutual Independence}

Events $(A,B,C)$ are \textbf{mutually independent} if
\[
\Pr(A,B)=\Pr(A)\Pr(B), \quad
\Pr(A,C)=\Pr(A)\Pr(C), \quad
\Pr(B,C)=\Pr(B)\Pr(C),
\]
and
\[
\Pr(A,B,C)=\Pr(A)\Pr(B)\Pr(C).
\]

\subsection{Communication Network Examples}

\begin{itemize}
  \item Probability of successful transmission is computed using independence
        of link availability.
  \item Modified networks were analyzed by enumerating disjoint successful paths.
\end{itemize}

\section{Types of Discrete Random Variables}

\subsection{Bernoulli Random Variable}

\textbf{Experiment:} Outcome is Success or Failure.

\textbf{Definition:}
\[
X =
\begin{cases}
1, & \text{Success} \\
0, & \text{Failure}
\end{cases}
\]

\textbf{PMF:}
\[
\Pr(X=1)=p, \quad \Pr(X=0)=1-p, \quad p\in(0,1).
\]

\textbf{Examples:}
\begin{itemize}
  \item Single coin toss
  \item Spam classification
\end{itemize}

\subsection{Binomial Random Variable}

\textbf{Experiment:} $n$ independent trials, each with success probability $p$.

\textbf{Random Variable:} $(X)$ = number of successes.

\textbf{Notation:} $X \sim B(n,p)$.

\textbf{PMF:}
\[
\Pr(X=i) = \binom{n}{i}p^i(1-p)^{n-i},
\quad i=0,1,\ldots,n.
\]

\subsection{Geometric Random Variable}

\textbf{Experiment:} Independent trials performed until first success.

\textbf{Random Variable:} $(X)$ = number of trials until success.

The PMF was derived in lecture.

\section{End Notes}

\begin{itemize}
  \item All results are restricted to material explicitly covered in Lecture 6.
  \item Proofs, derivations, and examples follow the lecture slides and board
        explanations exactly.
  \item This scribe is intended for \textbf{exam revision and reference}.
\end{itemize}

\end{document}
